\section{Introduction}\label{sec:intro}

Low-light video enhancement is ill-posed: the same dark clip admits many
plausible well-lit versions, and supervised models resolve the ambiguity with
statistics learned from one paired dataset. They do well on held-out clips of
that dataset and degrade on other sensors and scenes~\citep{bvi_rlv}. A prior
that does not come from the paired data is an obvious remedy, and video world
models are an attractive source: V-JEPA~\citep{vjepa,vjepa2} is trained to
predict the latent representation of masked spatio-temporal regions, so it is
usually described as modelling how scenes evolve rather than how they look.

That description is what makes the prior appealing, and this paper is about the
gap between it and what the prior is actually used for. We condition a
trainable Restormer-style pixel branch on frozen V-JEPA\,2 tokens through
zero-initialised residual cross-attention -- the simplest injection we could
write down -- and then spend the paper trying to establish that anything more
elaborate is warranted. It is not.

\paragraph{The prior works, and the gain is attributable.}
On held-out SDSD sequences the prior improves a matched network that differs
only in the tokens it reads by $\R{indomain@psnr}$\,dB PSNR-Y
$\R{indomain@psnr@ci}$ and $\R{indomain@lpips}$ LPIPS
$\R{indomain@lpips@ci}$, significant on both metrics across
\R{indomain@seeds} seeds with every seed agreeing in sign. Tokens from a
randomly initialised encoder of the same size are worth $\R{randomgap@psnr}$\,dB
less than pretrained ones $\R{randomgap@psnr@ci}$, so the gain comes from the
tokens' content and not from the extra conditioning path. The prior also
transfers: on each of three held-out domains it improves at least one metric
significantly with no significant degradation of the other
(Sec.~\ref{sec:exp:xdomain}).

\paragraph{The transfer result is invisible at the usual budget.}
Conditioning studies are commonly run at a few thousand iterations. At that
budget, on our data, one configuration's three seeds span 0.91\,dB while the
architectural effects being judged are 0.11--0.68\,dB, validation curves cross
between steps 1000 and 2000, and no run has plateaued. We re-ran every
architectural and cross-domain comparison at convergence and the verdicts moved:
of the three held-out domains that pass at convergence, none passed at the
short budget, and the one domain that was significant there is significant here
on a different metric. The short protocol also carries a bias against
additions, since anything with new parameters looks worse inside an
undertrained window -- which is the shape of every architectural improvement we
first believed we had.

\paragraph{What transfers is far less than the architecture assumes.}
Five ways of giving the pixel branch more of the representation -- positional
grounding, cross-frame keys, multi-layer keys, exposure decoupling, and a
learned reliability gate -- are at convergence statistically indistinguishable
from the simplest injection or significantly worse. Shuffling the input frames
changes the output by $\R{ordersens}$\,dB. Replacing all 256 spatial positions
of each tubelet with their mean is not significantly worse than the full field
($\R{suff@spatial@d}$\,dB, $\R{suff@spatial@ci}$), and neither is averaging
over time. Projecting each clip's 1024-dimensional tokens onto that clip's own
four principal directions does not hurt either ($\R{suff@rank4@d}$\,dB,
$\R{suff@rank4@ci}$) -- but four \emph{fixed} directions estimated on the
training split cost $\R{gsuff@global4}$\,dB, close to the
$\R{randomgap@psnr}$\,dB of discarding the tokens, and a branch trained on a
fixed 64-dimensional projection is significantly worse than one trained on the
full field ($\R{grank64@psnr}$\,dB, $\R{grank64@psnr@ci}$). Both halves are
needed to read the result correctly. What a converged branch relies on is a
low-rank subspace whose identity changes from clip to clip, so the reliance is
small while obtaining it still requires the encoder. The one design choice that
survives every test is depth: reading only the encoder's last block instead of
four blocks spanning its depth costs $\R{e2deep@psnr}$\,dB
($\R{e2deep@psnr@ci}$).

The temporal case is the sharpest, because the information is present and the
model declines it. Reordering a clip's frames changes V-JEPA's block-23 tokens
by 58\%, and low light does not degrade that structure. The converged model's
own latent dynamics are nonetheless almost orthogonal to the target's
(agreement $+0.034$ against $+0.968$ for a well-exposed clip under a mild
photometric change). Adding a loss that matches them raises the agreement to
$\R{h17@tdcos@cos}$ and makes perceptual quality and pixel-space temporal
consistency significantly \emph{worse} ($\R{h17@tdcos@lpips}$ LPIPS,
$\R{h17@tdcos@tlp}$ tLP gap). Latent temporal consistency and the kind the
metrics and the eye reward are not the same quantity.

\paragraph{Contributions.}
(i)~A controlled, converged attribution of a frozen video world model's benefit
for low-light video enhancement, significant in domain on both fidelity and
perceptual quality and significant on all three held-out domains, with
random-encoder and no-token controls.
(ii)~A measurement of \emph{what} a trained branch uses rather than how much it
gains: its reliance is low-rank, clip-adaptive, spatially near-uniform and
temporally near-constant, spread across the encoder's depth rather than its
channels, which accounts in one stroke for five architectural additions that do
not help -- and is distinguished, by training-time controls, from the stronger
and false claim that the representation could have been smaller.
(iii)~Evidence that forcing the unused temporal structure into the objective
engages the intended mechanism and degrades the result, with the instrument
that measures it and its calibrated ceiling.
(iv)~A protocol result: the training budget at which such comparisons are
usually made cannot resolve the effects being compared, demonstrated by
re-deciding six hypotheses at convergence.
All claims are pre-registered, and every number is generated from released
per-sequence results.
